\section{Discussão}
\label{sec:discussao}

\subsection{Pré-treino contra pós-treino}

As três questões ilustram dois regimes complementares de adaptação. O pré-treino
continuado (Questão 01) atua sobre a \emph{competência linguística} do modelo: sem
rótulos, apenas expondo-o ao corpus do domínio, reduziu a perplexidade em mais de
90\% e tornou suas gerações estilística e factualmente coerentes com o texto
administrativo do Piauí. Já o pós-treino (Questões 02 e 03) atua sobre o
\emph{comportamento} do modelo, ensinando-o a responder no formato
instrução/resposta. São objetivos distintos, e os resultados da Questão 01 mostram
que o ganho mais expressivo, em termos de modelagem de linguagem pura, vem da
exposição massiva ao domínio.

\subsection{SFT completo contra LoRA/QLoRA}

A comparação entre as Questões 02 e 03 é, em essência, uma análise de custo-benefício.
Os \textit{baselines} já coletados e o número de parâmetros treináveis
(Tabela~\ref{tab:comparativo}) deixam o trade-off evidente: o LoRA/QLoRA atualiza
0,44\% dos parâmetros e estimou treino cerca de 20 vezes mais rápido que o ajuste fino
completo no mesmo hardware. A hipótese a ser confirmada na execução final é a relatada
na literatura~\cite{hu2022lora,dettmers2023qlora}: que essa fração mínima de
parâmetros recupera a maior parte do ganho de qualidade do SFT completo, viabilizando
o ajuste em GPUs modestas.

\begin{table}[h]
  \centering
  \caption{Comparativo das três questões (valores ``depois'' de Q2 e Q3 pendentes de execução).}
  \label{tab:comparativo}
  \footnotesize
  \begin{tabularx}{\textwidth}{@{}l c c c c@{}}
    \toprule
    \textbf{Questão} & \textbf{Params. treinados} & \textbf{PPL antes} & \textbf{PPL depois} & \textbf{Acc. antes} \\
    \midrule
    01 -- Pré-treino   & 494M (100\%)   & 21,78 & \textbf{2,09} & 42,25\% \\
    02 -- SFT (full)   & 494M (100\%)   & 27,53 & \textit{em andamento} & 41,37\% \\
    03 -- LoRA/QLoRA   & 2,16M (0,44\%) & 34,51 & \textit{em andamento} & 38,42\% \\
    \bottomrule
  \end{tabularx}
\end{table}

\textbf{Ressalva de comparabilidade.} As perplexidades não são diretamente
comparáveis entre questões: a Questão 01 avalia sobre o DOMPI-2025, enquanto as
Questões 02 e 03 avaliam sobre o docentesDC; além disso, o \textit{baseline} da
Questão 03 é medido sobre o modelo quantizado em 4 bits. A comparação válida é sempre
\emph{antes contra depois dentro da mesma questão}.

\subsection{Limitações}

Cabe registrar as limitações do estudo. (i) As Questões 02 e 03 ainda não têm
métricas pós-treino, de modo que suas conclusões de qualidade são, por ora,
hipóteses fundamentadas no \textit{baseline} e na literatura, não em resultados
finais. (ii) A baixa perplexidade da Questão 01 é parcialmente atribuível à natureza
repetitiva do diário oficial, e não apenas à adaptação do modelo. (iii) Os pares de
QA gerados para o pós-treino têm qualidade desigual, conforme discutido na
Seção~\ref{sec:q2}, o que pode limitar o ganho do SFT antes das melhorias propostas.
(iv) Empregou-se um único modelo de 0,5~bilhão de parâmetros; o requisito opcional de
comparar modelos de tamanhos diferentes (p.\,ex.\ um Qwen2.5-1.5B via QLoRA) não foi
realizado.
